%% source: 01_intro (870w) -> target ~500w

\section{Introduction}\label{sec:intro}

Interactive vector maps on the web run two performance-critical paths.
The per-tile path fetches, decodes, and binds a client-side style to each tile's features at load time.
The per-frame path re-evaluates view-dependent expressions and submits draw calls to the GPU.
A slow network stalls the first, while a battery- or thermally-bound device stalls the second.
Libraries such as MapLibre GL~JS and Mapbox GL~JS deliver these maps by streaming tiled geometry to the browser and driving a WebGL or WebGPU rendering pipeline from a style document.
This architecture is what makes panning, zooming, and rotation smooth when it works, and stutter when it does not.

As deployments scale from personal projects to platforms serving millions, the cost of suboptimal performance compounds.
Larger tiles increase transfer time and bandwidth cost, redundant style layers waste rendering cycles, and unnecessary properties inflate decoding overhead.
On mobile devices these inefficiencies translate into battery drain and thermal throttling, and on unstable connections such as trains or aircraft they degrade the user experience directly.

Despite this, the map rendering community has largely treated style authoring and data preparation as independent concerns.
Styles are hand-tuned for visual appearance, while tiles are encoded with general-purpose content and compression.
The interaction between the two is left unexploited.
A style filter that matches on a property value may, after constant folding against tile-level statistics, simplify to a form that no longer references that property.
This in turn should allow an encoder to drop the corresponding column from every tile, a saving that neither style simplification nor data pruning could achieve alone.
We ask if these interactions yield synergestic or merely additive gains.

We investigate this automatic, style- and data-driven co-optimization resulting in these contributions:

\begin{itemize}
    \item A \textbf{formalization of joint style and tile co-optimization} as a transformation $\mathcal{O}: \mathcal{S} \times \mathcal{T} \to \mathcal{S} \times \mathcal{T}$ constrained by visual equivalence and strict idempotency.
    Prior work on compiler optimization and on columnar encoding has addressed one side in isolation, but neither has treated both jointly under a soundness criterion.
    \item A \textbf{three-level co-optimization strategy}.
    We first apply expression level peephole rewrites to a fixpoint.
    Then structural level runs (dead layer elimination, metadata refinement, and layer merging) pass over a style model.
    And finally, tile level optimisations employs \emph{style-driven tile shaving} (pruning unused properties/values/geometry types and feature ids) and an \emph{adaptive multi-strategy encoding} for the \ac{MLT} columnar format.
    \item An \textbf{end-to-end implementation} integrating the MapLibre style runtime and the \ac{MLT} tile pipeline.
    \item An \textbf{evaluation across four corpora} spanning 18 scenarios and a cumulative ablation, quantifying load-time, tile-volume, and frame-rate effects and resolving the additive-versus-synergistic question at both the aggregate and per-scenario levels.
\end{itemize}

Section~\ref{sec:background} gives a short overview of the background.
Section~\ref{sec:methods} builds on it to formalize the problem and detail the three levels.
Section~\ref{sec:results} reports the results, and Section~\ref{sec:discussion} interprets the findings.
